\section{Conclusion}
\label{sec:conclusion}

We have developed a substrate-consistency argument: a minimal deterministic hypothesis --- a 4D brane lattice embedded with codimension zero in a 4D Euclidean ambient --- together with five explicit mathematical bridges connecting it to the main structures any viable physical ontology must possess.
The ambient is fully symmetric; the timelike direction is selected by the Lorentzian sign structure of the brane action, and the gauge/gravity split is the inside observer's view after picking that direction.
The continuum model is stated as the long-wavelength effective description of the cubic substrate on a spacelike slice, using an isotropic
hyperelastic action built from the induced metric and a homogeneous prestretch parameter $\alpha$ as the analytic
baseline. The bare cubic stencil must still pass the separate dispersion/isotropy tests, because the Cauchy relation
alone does not guarantee an isotropic leading-order acoustic tensor.

A central clarification of the present revision is that the essential nonlinearity is already geometric.
Even when the microscopic link energy is Hooke-linear in scalar extension, the full Euclidean link length in $\R^4$
and the induced metric generate nonlinear amplitude--lateral coupling.
That same mechanism is intended to underwrite contraction, self-guidance, and localization.

Applying the complex-envelope description per band-isolated wavepacket, we constructed Berry and Wilczek--Zee
connections from adiabatic eigenbundle transport and interpreted them as the natural gauge variables of the
per-wavepacket coarse-grained envelope description.
On a cubic lattice, an axis-aligned triplet sector provides the natural candidate internal arena in which a non-Abelian
transport structure can emerge.
Localized particle-like states were then posed as self-adjoint eigenproblems around symmetric backgrounds, with
confinement attributed to geometric self-guidance rather than to fundamental point charges or externally imposed
potentials.

The most important consequence is methodological.
The theory is now stated in a way that makes its next decisive test clear: the search for stable, localized,
three-component baryon-like modes on the cubic substrate.
Such a calculation would simultaneously probe the geometric confinement mechanism, the axis-aligned triplet structure,
and the viability of coarse-grained isotropy.
Whether that program succeeds or fails is more important than further verbal extrapolation.
The immediate task is therefore to turn the consolidated theory into direct computation.

The interpretive commitment that closes the present paper --- the retrocausal worldtube reading of \Cref{sec:bell-interpretation} --- is not a free embellishment.
It is forced on any local + deterministic substrate by Bell's theorem, and it shapes how the simulation programme should ultimately read its own results: matter and antimatter as forward- and backward-propagating soliton worldtubes, and entanglement as continuity of one 4D object that branches at a V-vertex in the past.
